\documentclass[11pt]{article}

\usepackage[a4paper,margin=1in]{geometry}
\usepackage[T1]{fontenc}
\usepackage[utf8]{inputenc}
\usepackage{lmodern}
\usepackage{microtype}
\usepackage{setspace}
\usepackage{parskip}
\usepackage{hyperref}
\usepackage{enumitem}
\usepackage{longtable}
\usepackage{array}
\usepackage{booktabs}

\onehalfspacing

\title{Minimal Reviewable Evidence Specification\\
for Constitutional AI}
\author{Panagiotis Kalomoirakis\\
Independent Researcher, Synkyria Project}
\date{April 2026}

\begin{document}

\maketitle

\begin{center}
\small \textcopyright\ Panagiotis Kalomoirakis 2026. All rights reserved.
\end{center}

\begin{abstract}
This document specifies a minimal reviewable evidence surface for AI systems claiming constitutional operation under finite capacity. It does not define a full audit standard, implementation protocol, or legal compliance framework. Its narrower purpose is to state the minimum review artefacts required for a constitutional claim to be inspectable rather than merely asserted. The specification defines the evidence objects, their minimum fields, integrity requirements, scope declarations, certificate discipline, and non-claim boundaries. It is intended to function as a companion object to a constitutional article draft and one or more application profiles.
\end{abstract}

\section{Status of this document}

This document is a \textbf{minimal evidence specification}. It does not prove that any AI system is safe, ethical, lawful, unbiased, or fit for deployment in a specific domain. It specifies only the minimum evidence conditions under which a constitutional claim may become reviewable.

Its function is therefore narrow:

\begin{quote}
to prevent constitutional language from floating free of inspectable evidence.
\end{quote}

\section{Purpose}

A constitutional AI claim must not remain a slogan. If a system claims that it:
\begin{itemize}[leftmargin=2em]
    \item distinguishes admissibility from execution,
    \item can hold or refuse lawfully,
    \item preserves witness,
    \item remains reviewable under finite capacity,
\end{itemize}
then a minimum evidence layer must exist through which those claims may be inspected.

This specification defines that minimum layer.

\section{Scope}

This specification applies to AI systems that make claims such as:
\begin{itemize}[leftmargin=2em]
    \item constitutional operation,
    \item reviewable refusal,
    \item explicit holding,
    \item witness-preserving boundary governance,
    \item bounded admissibility discipline,
    \item or certificate-bearing accountable operation.
\end{itemize}

It may be used together with:
\begin{itemize}[leftmargin=2em]
    \item a constitutional article draft,
    \item an application profile,
    \item a domain-specific deployment profile,
    \item or a runtime demonstration pack.
\end{itemize}

\section{Evidence principle}

The core evidence principle is:

\begin{quote}
No constitutional act without witness; no witness without scope; no scope without review conditions.
\end{quote}

Accordingly, a minimal reviewable evidence package must always make visible:
\begin{itemize}[leftmargin=2em]
    \item \textbf{what} act occurred,
    \item \textbf{why} it was declared,
    \item \textbf{under which system and profile} it occurred,
    \item \textbf{what the evidence does and does not support},
    \item and \textbf{how integrity is preserved}.
\end{itemize}

\section{Minimal evidence package}

A minimal reviewable evidence package should contain, at minimum, the following artefacts:

\begin{enumerate}[leftmargin=2em,label=\arabic*.]
    \item a \texttt{manifest.json}
    \item an action trace such as \texttt{actions.jsonl}
    \item optionally, a frame or state trace such as \texttt{frames.jsonl}
    \item an integrity file such as \texttt{checksums.sha256}
    \item a short natural-language scope note such as \texttt{README.md}
    \item and, where a verdict is issued, a \texttt{certificate.json} or certificate envelope
\end{enumerate}

These filenames are illustrative rather than mandatory. What is mandatory is the review function they serve.

\section{Required artefact functions}

\subsection{Manifest}

The manifest identifies the evidence package.

It should minimally declare:
\begin{itemize}[leftmargin=2em]
    \item system name,
    \item profile name,
    \item version,
    \item date or timestamp,
    \item scenario or run label,
    \item applicable domain profile,
    \item relevant configuration scope,
    \item and package contents.
\end{itemize}

\subsection{Action trace}

The action trace records constitutionally significant acts.

It should minimally allow a reviewer to inspect:
\begin{itemize}[leftmargin=2em]
    \item event order,
    \item the declared act,
    \item the associated reason label or reason family,
    \item the interaction or step identifier,
    \item and any scoped witness object or pointer attached to the act.
\end{itemize}

\subsection{Frame or state trace}

Where needed, the frame trace records the surrounding runtime surface within which an act occurred.

It may include:
\begin{itemize}[leftmargin=2em]
    \item local telemetry,
    \item bounded pressure indicators,
    \item tool state,
    \item review-relevant flags,
    \item or other declared witness fields.
\end{itemize}

It is not mandatory in all settings, but where omitted, the package must still remain reviewable through the available artefacts.

\subsection{Integrity file}

The integrity file allows the reviewer to verify that the package has not been silently altered.

It should minimally support file-level integrity checking across the included artefacts.

\subsection{Scope note}

The scope note explains:
\begin{itemize}[leftmargin=2em]
    \item what is being claimed,
    \item what is not being claimed,
    \item what profile or scenario the package belongs to,
    \item what limitations govern interpretation,
    \item and what a reviewer may legitimately conclude from the package.
\end{itemize}

\subsection{Certificate object}

Where a certificate is issued, it must remain subordinate to the evidence.

A certificate must not replace the evidence package. It may summarise, but it must not erase ambiguity, scope limits, or lower-layer qualifications.

\section{Minimum fields}

\subsection{Minimum fields for \texttt{manifest.json}}

A minimal manifest should include fields equivalent to the following:

\begin{center}
\begin{tabular}{>{\ttfamily}p{4cm} p{9cm}}
\toprule
field & function \\
\midrule
system\_name & identifies the system under review \\
profile\_name & identifies the relevant application or domain profile \\
version & identifies the system/profile version \\
run\_label & identifies the specific run or evidence instance \\
timestamp\_utc & anchors the package in time \\
package\_contents & lists included artefacts \\
scope & states what this package is meant to evidence \\
non\_claims & states what this package does not establish \\
integrity\_method & identifies how integrity is checked \\
\bottomrule
\end{tabular}
\end{center}

\subsection{Minimum fields for an action trace}

Each constitutionally significant line or event should minimally expose fields equivalent to:

\begin{center}
\begin{tabular}{>{\ttfamily}p{4cm} p{9cm}}
\toprule
field & function \\
\midrule
step / event\_id & preserves event order \\
timestamp\_utc & anchors the event in time \\
action\_type & records the declared act \\
reason\_code & records why the act was declared \\
witness & attaches a bounded review object or pointer \\
profile\_version & ties the act to a declared profile/version \\
\bottomrule
\end{tabular}
\end{center}

\subsection{Minimum fields for a certificate}

If used, a certificate should minimally expose fields equivalent to:

\begin{center}
\begin{tabular}{>{\ttfamily}p{4cm} p{9cm}}
\toprule
field & function \\
\midrule
certificate\_type & states the type of verdict object \\
scope & states what is being certified \\
status & states the verdict or classification \\
qualifications & states lower-layer limits or caveats \\
evidence\_refs & points back to the supporting evidence \\
non\_claims & prevents over-interpretation \\
\bottomrule
\end{tabular}
\end{center}

\section{Act evidence requirements}

\subsection{\texttt{ACCEPT}}

An \texttt{ACCEPT} event should be reviewable as an explicit admission of continuation.

Minimal review requirement:
\begin{itemize}[leftmargin=2em]
    \item the act must be explicit,
    \item it must not be inferred merely from the presence of output,
    \item and it must remain tied to a declared profile and reasonable scope.
\end{itemize}

\subsection{\texttt{HOLD}}

A \texttt{HOLD} event must not appear as passive delay, silent wait, or hidden buffering.

Minimal review requirement:
\begin{itemize}[leftmargin=2em]
    \item the act must be explicit,
    \item it must carry a reason family,
    \item and the package must make visible that lawful continuation was deferred rather than silently abandoned.
\end{itemize}

\subsection{\texttt{REFUSE}}

A \texttt{REFUSE} event must not be substituted by timeout, unexplained absence, or opaque failure.

Minimal review requirement:
\begin{itemize}[leftmargin=2em]
    \item the refusal must be explicit,
    \item it must carry a reason family,
    \item and it must remain inspectable as a constitutional act rather than a transport or infrastructure accident.
\end{itemize}

\subsection{Extension acts}

If the system declares extension acts such as \texttt{ESCALATE}, \texttt{TRANSFER}, or \texttt{RE-DESCRIBE}, the evidence package must make visible:
\begin{itemize}[leftmargin=2em]
    \item that the extension act occurred,
    \item under what declared scope it is valid,
    \item and how responsibility and witness are preserved across the transition.
\end{itemize}

\section{Reason-label discipline}

Reason labels shall be:
\begin{itemize}[leftmargin=2em]
    \item declared,
    \item scoped,
    \item semantically stable across a package or profile,
    \item and sufficient for later review.
\end{itemize}

Reason labels need not expose raw thresholds, coefficients, or proprietary mechanisms. However, they must not be so vague that the act becomes unreviewable.

Illustrative reason families include:
\begin{itemize}[leftmargin=2em]
    \item \texttt{OVERLOAD}
    \item \texttt{UNDER\_REVIEW}
    \item \texttt{SAFETY\_BOUNDARY}
    \item \texttt{SCOPE\_LIMIT}
    \item \texttt{INSUFFICIENT\_BASIS}
    \item \texttt{ESCALATION\_REQUIRED}
\end{itemize}

\section{Witness discipline}

A witness object may be minimal, but it must not be empty.

At minimum, the witness must help the reviewer answer:

\begin{quote}
Why should this event count as a constitutionally meaningful act under the declared profile?
\end{quote}

A witness may be:
\begin{itemize}[leftmargin=2em]
    \item inline in the action trace,
    \item attached as a bounded object,
    \item or referenced via a linked artefact,
\end{itemize}
provided that it remains accessible and scoped.

\section{Scope and non-claim discipline}

Every evidence package must explicitly distinguish:
\begin{itemize}[leftmargin=2em]
    \item \textbf{what it supports},
    \item \textbf{what it illustrates},
    \item \textbf{what it does not prove},
    \item and \textbf{under what profile or scenario it is valid}.
\end{itemize}

This is necessary to prevent:
\begin{itemize}[leftmargin=2em]
    \item witness inflation,
    \item certificate overreach,
    \item demo-to-universal drift,
    \item and public-facing overclaim.
\end{itemize}

\section{Qualification inheritance}

No certificate or higher-level summary may erase unresolved lower-layer weakness.

Therefore:
\begin{itemize}[leftmargin=2em]
    \item if the evidence is scenario-bound, the certificate must say so;
    \item if the witness is partial, the certificate must say so;
    \item if lower-layer instability remains unresolved, the higher verdict must inherit that qualification explicitly.
\end{itemize}

A reviewer must never be asked to treat a high-level verdict as stronger than the evidence beneath it.

\section{Protected non-disclosure}

This specification does not require full exposure of:
\begin{itemize}[leftmargin=2em]
    \item raw thresholds,
    \item model weights,
    \item proprietary parameters,
    \item or security-sensitive internals.
\end{itemize}

But if such material is withheld, the package must still preserve sufficient witness for meaningful review.

Non-disclosure is constitutionally valid only if reviewability survives.

\section{Minimal package example}

A minimal reviewable evidence package may therefore look like:

\begin{verbatim}
/
|- manifest.json
|- actions.jsonl
|- frames.jsonl           (optional but recommended where relevant)
|- checksums.sha256
|- certificate.json       (optional, if a verdict is issued)
|- README.md
\end{verbatim}

The exact filenames may vary. The constitutional requirement lies in the review function, not in the filename itself.

\section{Failure modes this specification is meant to prevent}

This specification is designed to reduce the risk of:
\begin{itemize}[leftmargin=2em]
    \item silent degradation,
    \item timeout-as-refusal,
    \item unscoped refusal,
    \item witness-free continuation,
    \item certificate theatre,
    \item opaque escalation,
    \item non-reviewable holding,
    \item and scope drift between evidence and public claim.
\end{itemize}

\section{Non-claims}

A package satisfying this specification does \emph{not} by itself prove:
\begin{itemize}[leftmargin=2em]
    \item ethical legitimacy,
    \item legal compliance,
    \item universal safety,
    \item fairness,
    \item domain fitness,
    \item or deployment readiness.
\end{itemize}

It proves something narrower:

\begin{quote}
that the constitutional claim being made has been attached to a minimum reviewable evidence surface.
\end{quote}

\section{Minimal conformance checklist}

A package minimally conforms to this specification only if all of the following are true:

\begin{enumerate}[leftmargin=2em,label=\arabic*.]
    \item It identifies the system, profile, version, and evidence instance.
    \item It preserves an explicit action trace for constitutionally significant acts.
    \item It includes a manifest stating scope and non-claims.
    \item It includes an integrity mechanism.
    \item It provides reason-label discipline for non-trivial acts.
    \item It preserves bounded witness for later review.
    \item It distinguishes evidence from certificate.
    \item It preserves qualification inheritance across levels.
    \item It does not allow non-disclosure to erase reviewability.
    \item It does not present the package as stronger than its declared scope.
\end{enumerate}

\section{Closing statement}

A constitutional claim without evidence is rhetoric.  
Evidence without scope is ambiguity.  
A certificate without qualification is overreach.  

This specification therefore establishes a minimal rule: if an AI system claims constitutional operation under finite capacity, it must leave behind a reviewable evidence surface through which its boundary acts can be inspected. Without that surface, constitutional language remains structurally weak.

\end{document}